\documentclass{amsart}
% For dealing with references we use the comment environment
\usepackage{verbatim}
\newenvironment{reference}{\comment}{\endcomment}
%\newenvironment{reference}{}{}
\newenvironment{slogan}{\comment}{\endcomment}
\newenvironment{history}{\comment}{\endcomment}

% For commutative diagrams we use Xy-pic
\usepackage[all]{xy}

% We use 2cell for 2-commutative diagrams.
\xyoption{2cell}
\UseAllTwocells

% We use multicol for the list of chapters between chapters
\usepackage{multicol}

% This is generally recommended for better output
\usepackage{lmodern}
\usepackage[T1]{fontenc}

% For cross-file-references

% Package for hypertext links:
\usepackage{hyperref}

% For any local file, say "hello.tex" you want to link to please

% Theorem environments.
%
\theoremstyle{plain}
\newtheorem{theorem}[subsection]{Theorem}
\newtheorem{proposition}[subsection]{Proposition}
\newtheorem{lemma}[subsection]{Lemma}

\theoremstyle{definition}
\newtheorem{definition}[subsection]{Definition}
\newtheorem{example}[subsection]{Example}
\newtheorem{exercise}[subsection]{Exercise}
\newtheorem{situation}[subsection]{Situation}

\theoremstyle{remark}
\newtheorem{remark}[subsection]{Remark}
\newtheorem{remarks}[subsection]{Remarks}

\numberwithin{equation}{subsection}

% Macros
%
\def\lim{\mathop{\mathrm{lim}}\nolimits}
\def\colim{\mathop{\mathrm{colim}}\nolimits}
\def\Spec{\mathop{\mathrm{Spec}}}
\def\Hom{\mathop{\mathrm{Hom}}\nolimits}
\def\Ext{\mathop{\mathrm{Ext}}\nolimits}
\def\SheafHom{\mathop{\mathcal{H}\!\mathit{om}}\nolimits}
\def\SheafExt{\mathop{\mathcal{E}\!\mathit{xt}}\nolimits}
\def\Sch{\mathit{Sch}}
\def\Mor{\mathop{\mathrm{Mor}}\nolimits}
\def\Ob{\mathop{\mathrm{Ob}}\nolimits}
\def\Sh{\mathop{\mathit{Sh}}\nolimits}
\def\NL{\mathop{N\!L}\nolimits}
\def\CH{\mathop{\mathrm{CH}}\nolimits}
\def\proetale{{pro\text{-}\acute{e}tale}}
\def\etale{{\acute{e}tale}}
\def\QCoh{\mathit{QCoh}}
\def\Ker{\mathop{\mathrm{Ker}}}
\def\Im{\mathop{\mathrm{Im}}}
\def\Coker{\mathop{\mathrm{Coker}}}
\def\Coim{\mathop{\mathrm{Coim}}}

% Boxtimes
%
\DeclareMathSymbol{\boxtimes}{\mathbin}{AMSa}{"02}

%
% Macros for moduli stacks/spaces
%
\def\QCohstack{\mathcal{QC}\!\mathit{oh}}
\def\Cohstack{\mathcal{C}\!\mathit{oh}}
\def\Spacesstack{\mathcal{S}\!\mathit{paces}}
\def\Quotfunctor{\mathrm{Quot}}
\def\Hilbfunctor{\mathrm{Hilb}}
\def\Curvesstack{\mathcal{C}\!\mathit{urves}}
\def\Polarizedstack{\mathcal{P}\!\mathit{olarized}}
\def\Complexesstack{\mathcal{C}\!\mathit{omplexes}}
% \Pic is the operator that assigns to X its picard group, usage \Pic(X)
% \Picardstack_{X/B} denotes the Picard stack of X over B
% \Picardfunctor_{X/B} denotes the Picard functor of X over B
\def\Pic{\mathop{\mathrm{Pic}}\nolimits}
\def\Picardstack{\mathcal{P}\!\mathit{ic}}
\def\Picardfunctor{\mathrm{Pic}}
\def\Deformationcategory{\mathcal{D}\!\mathit{ef}}

\setlength{\emergencystretch}{3em}
\begin{document}
\title[Stacks Project, Tag 053R]{482 --- Vertical locally principal subschemes occupy fibre components under the irreducible dominant finite-type and nonincreasing fibre-dimension hypotheses}
\maketitle
\noindent Copyright (C) 2005--2025 Johan de Jong. Stacks Project authors. Source: \href{https://stacks.math.columbia.edu/tag/053R}{Tag 053R}.

\noindent Editorial difficulty note (not author text). Difficulty H2: The Noetherian argument combines the principal ideal theorem with specialization chains and the fibre-dimension bound. The general case requires a careful finite-type integer-algebra approximation preserving generic invertibility, both fibre dimensions and avoidance of all minimal primes, so that the same contradiction descends..

\noindent Editorial provenance note (not author text). History: excerpt from revision \texttt{a0444\allowbreak 6e57e\allowbreak c1fbc\allowbreak 252a8\allowbreak 71afc\allowbreak ec775\allowbreak 2fb28\allowbreak 07b14}, more-morphisms.tex, lines 14537--14694. Reference presentation was adapted; original-author context and core support blocks were retained or added. The mathematical wording is unchanged. Licensed under GNU FDL 1.2 or later; no invariant sections or cover texts. See ../licenses/COPYING. Context blocks reproduce author definitions and supporting results; they are not additional counted problems.

\begin{lemma}
\label{lemma-locally-principal-vertical}
Let $f : X \to S$ be a morphism of schemes.
Let $Z \subset X$ be a closed subscheme.
Let $s \in S$.
Assume
\begin{enumerate}
\item $S$ is irreducible with generic point $\eta$,
\item $X$ is irreducible,
\item $f$ is dominant,
\item $f$ is locally of finite type,
\item $\dim(X_s) \leq \dim(X_\eta)$,
\item $Z$ is locally principal in $X$, and
\item $Z_\eta = \emptyset$.
\end{enumerate}
Then the fibre $Z_s$ is (set theoretically) a union of
irreducible components of $X_s$.
\end{lemma}

\begin{proof}
Let $X_{red}$ denote the reduction of $X$. Then $Z \cap X_{red}$ is
a locally principal closed subscheme of $X_{red}$, see
Divisors, Lemma \href{https://stacks.math.columbia.edu/tag/053P}{lemma pullback locally principal (Tag 053P)}.
Hence we may assume that $X$ is reduced. In other words $X$ is integral, see
Properties, Lemma \href{https://stacks.math.columbia.edu/tag/01ON}{lemma characterize integral (Tag 01ON)}.
In this case the morphism $X \to S$ factors through $S_{red}$, see
Schemes, Lemma \href{https://stacks.math.columbia.edu/tag/0356}{lemma map into reduction (Tag 0356)}.
Thus we may replace $S$ by $S_{red}$ and assume that $S$ is integral too.

\medskip\noindent
The assertion that $f$ is dominant signifies that the generic point of $X$
is mapped to $\eta$, see
Morphisms,
Lemma \href{https://stacks.math.columbia.edu/tag/01RM}{lemma dominant finite number irreducible components (Tag 01RM)}.
Moreover, the scheme $X_\eta$ is an integral scheme which is locally of
finite type over the field $\kappa(\eta)$. Hence
$d = \dim(X_\eta) \geq 0$ is equal to $\dim_\xi(X_\eta)$ for
every point $\xi$ of $X_\eta$, see
Algebra, Lemmas \href{https://stacks.math.columbia.edu/tag/00OS}{lemma dimension spell it out (Tag 00OS)} and
\href{https://stacks.math.columbia.edu/tag/00OT}{lemma dimension at a point finite type over field (Tag 00OT)}.
In view of
Morphisms, Lemma \href{https://stacks.math.columbia.edu/tag/02FZ}{lemma openness bounded dimension fibres (Tag 02FZ)}
and condition (5) we conclude that $\dim_x(X_s) = d$
for every $x \in X_s$.

\medskip\noindent
In the Noetherian case the assertion can be proved as follows.
If the lemma does not holds there exists $x \in Z_s$ which is a generic
point of an irreducible component of $Z_s$ but not a generic point
of any irreducible component of $X_s$. Then we see that
$\dim_x(Z_s) \leq d - 1$, because $\dim_x(X_s) = d$ and in a neighbourhood
of $x$ in $X_s$ the closed subscheme $Z_s$ does not contain any of the
irreducible components of $X_s$. Hence after replacing $X$ by an
open neighbourhood of $x$ we may assume that
$\dim_z(Z_{f(z)}) \leq d - 1$ for all $z \in Z$, see
Morphisms, Lemma \href{https://stacks.math.columbia.edu/tag/02FZ}{lemma openness bounded dimension fibres (Tag 02FZ)}.
Let $\xi' \in Z$ be a generic point of an irreducible component of $Z$
and set $s' = f(\xi)$. As $Z \not = X$ is locally principal we see that
$\dim(\mathcal{O}_{X, \xi}) = 1$, see
Algebra, Lemma \href{https://stacks.math.columbia.edu/tag/00KV}{lemma minimal over 1 (Tag 00KV)}
(this is where we use $X$ is Noetherian).
Let $\xi \in X$ be the generic point of $X$ and
let $\xi_1$ be a generic point of any irreducible component
of $X_{s'}$ which contains $\xi'$. Then we see that we have
the specializations
$$
\xi \leadsto \xi_1 \leadsto \xi'.
$$
As $\dim(\mathcal{O}_{X, \xi}) = 1$ one of the two specializations
has to be an equality.
By assumption $s' \not = \eta$, hence the first specialization
is not an equality.
Hence $\xi' = \xi_1$ is a generic point of an irreducible component of
$X_{s'}$. Applying
Morphisms, Lemma \href{https://stacks.math.columbia.edu/tag/02FZ}{lemma openness bounded dimension fibres (Tag 02FZ)}
one more time this implies
$\dim_{\xi'}(Z_{s'}) = \dim_{\xi'}(X_{s'}) \geq \dim(X_\eta) = d$
which gives the desired contradiction.

\medskip\noindent
In the general case we reduce to the Noetherian case as follows.
If the lemma is false then there exists a point
$x \in X$ lying over $s$ such that $x$ is a generic point of an
irreducible component of $Z_s$, but
not a generic point of any of the irreducible components of $X_s$.
Let $U \subset S$ be an affine neighbourhood of $s$ and let
$V \subset X$ be an affine neighbourhood of $x$ with $f(V) \subset U$.
Write $U = \Spec(A)$ and $V = \Spec(B)$ so that $f|_V$
is given by a ring map $A \to B$. Let $\mathfrak q \subset B$,
resp.\ $\mathfrak p \subset A$ be the prime corresponding to $x$, resp.\ $s$.
After possibly shrinking $V$ we may assume $Z \cap V$ is cut out by
some element $g \in B$. Denote $K$ the fraction field of $A$.
What we know at this point is the following:
\begin{enumerate}
\item $A \subset B$ is a finitely generated extension of domains,
\item the element $g \otimes 1$ is invertible in $B \otimes_A K$,
\item $d = \dim(B \otimes_A K) = \dim(B \otimes_A \kappa(\mathfrak p))$,
\item $g \otimes 1$ is not a unit of $B \otimes_A \kappa(\mathfrak p)$, and
\item $g \otimes 1$ is not in any of the minimal primes of
$B \otimes_A \kappa(\mathfrak p)$.
\end{enumerate}
We are seeking a contradiction.

\medskip\noindent
Pick elements $x_1, \ldots, x_n \in B$ which generate $B$ over $A$.
For a finitely generated $\mathbf{Z}$-algebra $A_0 \subset A$
let $B_0 \subset B$ be the $A_0$-subalgebra generated by
$x_1, \ldots, x_n$, denote $K_0$ the fraction field of $A_0$, and set
$\mathfrak p_0 = A_0 \cap \mathfrak p$.
We claim that when $A_0$ is large enough then (1) -- (5) also hold for
the system $(A_0 \subset B_0, g, \mathfrak p_0)$.

\medskip\noindent
We prove each of the conditions in turn. Part (1) holds by construction.
For part (2) write $(g \otimes 1) h = 1$ for some
$h \otimes 1/a \in B \otimes_A K$. Write
$g = \sum a_I x^I$, $h = \sum a'_I x^I$ (multi-index notation)
for some coefficients $a_I, a'_I \in A$. As soon as $A_0$ contains
$a$ and the $a_I, a'_I$ then (2) holds because
$B_0 \otimes_{A_0} K_0 \subset B \otimes_A K$ (as localizations of the
injective map $B_0 \to B$).
To achieve (3) consider the exact sequence
$$
0 \to I \to A[X_1, \ldots, X_n] \to B \to 0
$$
which defines $I$ where the second map sends $X_i$ to $x_i$. Since $\otimes$
is right exact we see that $I \otimes_A K$, respectively
$I \otimes_A \kappa(\mathfrak p)$ is the kernel of the surjection
$K[X_1, \ldots, X_n] \to B \otimes_A K$, respectively
$\kappa(\mathfrak p)[X_1, \ldots, X_n] \to B \otimes_A \kappa(\mathfrak p)$.
As a polynomial ring over a field is Noetherian
there exist finitely many elements $h_j \in I$, $j = 1, \ldots, m$
which generate $I \otimes_A K$ and $I \otimes_A \kappa(\mathfrak p)$.
Write $h_j = \sum a_{j, I}X^I$. As soon as
$A_0$ contains all $a_{j, I}$ we get to the situation where
$$
B_0 \otimes_{A_0} K_0 \otimes_{K_0} K = B \otimes_A K
\quad\text{and}\quad
B_0 \otimes_{A_0} \kappa(\mathfrak p_0)
\otimes_{\kappa(\mathfrak p_0)} \kappa(\mathfrak p)
=
B \otimes_A \kappa(\mathfrak p).
$$
By either
Morphisms, Lemma \href{https://stacks.math.columbia.edu/tag/02FY}{lemma dimension fibre after base change (Tag 02FY)}
or
Algebra, Lemma \href{https://stacks.math.columbia.edu/tag/00P3}{lemma dimension preserved field extension (Tag 00P3)}
we see that the dimension equalities of (3) are satisfied.
Part (4) is immediate. As
$B_0 \otimes_{A_0} \kappa(\mathfrak p_0) \subset
B \otimes_A \kappa(\mathfrak p)$ each minimal prime of
$B_0 \otimes_{A_0} \kappa(\mathfrak p_0)$ lies under a minimal
prime of $B \otimes_A \kappa(\mathfrak p)$ by
Algebra, Lemma \href{https://stacks.math.columbia.edu/tag/00FL}{lemma image dense generic points (Tag 00FL)}.
This implies that (5) holds.
In this way we reduce the problem to the Noetherian case which we
have dealt with above.
\end{proof}
\end{document}
